% ============================================================================
% preamble.tex — 包、字体、排版、数学环境
% ============================================================================

% --- 中文支持 (XeLaTeX) ---
\usepackage{ctex}              % 中文排版核心
\usepackage{fontspec}          % 字体选择

% --- 页面布局 ---
\usepackage[
  a4paper,
  top=2.5cm, bottom=2.5cm,
  left=2.5cm, right=2.5cm
]{geometry}

% --- 数学 ---
\usepackage{amsmath, amssymb, amsthm}
\usepackage{mathrsfs}          % \mathscr 花体
\usepackage{mathtools}         % dcases, coloneqq 等
\usepackage{bm}                % 粗体数学符号

% --- 图表 ---
\usepackage{graphicx}
\usepackage{float}
\usepackage{booktabs}
\usepackage{longtable}
\usepackage{tabularx}

% --- 绘图 ---
\usepackage{tikz}
\usetikzlibrary{
  positioning, calc, arrows.meta,
  decorations.pathreplacing,
  patterns,
}

% --- 交叉引用与超链接 ---
\usepackage{hyperref}
\hypersetup{
  colorlinks=true,
  linkcolor=blue!60!black,
  citecolor=green!40!black,
  urlcolor=cyan!50!black,
}
\usepackage[nameinlink]{cleveref}  % 智能引用: \cref

% --- 索引与参考文献 ---
\usepackage{makeidx}
\makeindex

% --- 颜色 ---
\usepackage{xcolor}

% --- 彩色定理框 ---
\usepackage[most]{tcolorbox}
\tcbuselibrary{breakable, skins}

% ============================================================================
% 定理环境（amsthm 定义 + tcolorbox 包装，不同类型不同颜色）
% ============================================================================

% --- 先用 amsthm 定义环境 ---
\theoremstyle{plain}
\newtheorem{theorem}{定理}[chapter]
\newtheorem{proposition}[theorem]{命题}
\newtheorem{lemma}[theorem]{引理}
\newtheorem{corollary}[theorem]{推论}
\newtheorem{conjecture}[theorem]{猜想}

\theoremstyle{definition}
\newtheorem{definition}[theorem]{定义}
\newtheorem{example}[theorem]{例}
\newtheorem{exercise}{习题}[chapter]

\theoremstyle{remark}
\newtheorem{remark}[theorem]{注记}
\newtheorem{notation}[theorem]{记号}

% --- 用 tcolorbox 给每类环境加彩色边框 ---
% 定理/命题/引理/推论/猜想：蓝色左边条
\tcolorboxenvironment{theorem}{
  blanker, breakable, left=5mm,
  before skip=10pt, after skip=10pt,
  borderline west={2pt}{0pt}{blue!60!black},
  colback=blue!3,
}
\tcolorboxenvironment{proposition}{
  blanker, breakable, left=5mm,
  before skip=10pt, after skip=10pt,
  borderline west={2pt}{0pt}{blue!50!black},
  colback=blue!3,
}
\tcolorboxenvironment{lemma}{
  blanker, breakable, left=5mm,
  before skip=10pt, after skip=10pt,
  borderline west={2pt}{0pt}{blue!50!black},
  colback=blue!3,
}
\tcolorboxenvironment{corollary}{
  blanker, breakable, left=5mm,
  before skip=10pt, after skip=10pt,
  borderline west={2pt}{0pt}{blue!50!black},
  colback=blue!3,
}
\tcolorboxenvironment{conjecture}{
  blanker, breakable, left=5mm,
  before skip=10pt, after skip=10pt,
  borderline west={2pt}{0pt}{blue!50!black},
  colback=blue!3,
}

% 定义：绿色左边条
\tcolorboxenvironment{definition}{
  blanker, breakable, left=5mm,
  before skip=10pt, after skip=10pt,
  borderline west={2pt}{0pt}{green!50!black},
  colback=green!3,
}

% 例：橙色左边条
\tcolorboxenvironment{example}{
  blanker, breakable, left=5mm,
  before skip=10pt, after skip=10pt,
  borderline west={2pt}{0pt}{orange!70!black},
  colback=orange!3,
}

% 注记：灰色左边条
\tcolorboxenvironment{remark}{
  blanker, breakable, left=5mm,
  before skip=10pt, after skip=10pt,
  borderline west={2pt}{0pt}{gray!50},
  colback=gray!5,
}

% 习题：紫色左边条
\tcolorboxenvironment{exercise}{
  blanker, breakable, left=5mm,
  before skip=10pt, after skip=10pt,
  borderline west={2pt}{0pt}{violet!60!black},
  colback=violet!3,
}

% ============================================================================
% 常用数学命令
% ============================================================================

% 数域
\newcommand{\N}{\mathbb{N}}
\newcommand{\Z}{\mathbb{Z}}
\newcommand{\Q}{\mathbb{Q}}
\newcommand{\R}{\mathbb{R}}
\newcommand{\C}{\mathbb{C}}
\newcommand{\F}{\mathbb{F}}          % 有限域

% 上半平面
\newcommand{\HH}{\mathbb{H}}         % 上半平面 H

% 模群与矩阵群
\newcommand{\SL}{\mathrm{SL}}
\newcommand{\GL}{\mathrm{GL}}
\newcommand{\PSL}{\mathrm{PSL}}
\newcommand{\PGL}{\mathrm{PGL}}
\newcommand{\Sp}{\mathrm{Sp}}
\newcommand{\Mat}{\mathrm{Mat}}

% 同余子群
\newcommand{\GammaO}{\Gamma_0}       % Γ₀(N)
\newcommand{\GammaI}{\Gamma_1}       % Γ₁(N)

% 模形式空间
\newcommand{\Mk}{\mathcal{M}}        % M_k
\newcommand{\Sk}{\mathcal{S}}        % S_k (尖点形式)

% Hecke 算子
\newcommand{\Tk}{T}                   % T_n

% L-函数
\newcommand{\Lfun}{L}

% 自守表示
\newcommand{\Aut}{\mathrm{Aut}}
\newcommand{\Gal}{\mathrm{Gal}}
\newcommand{\Frob}{\mathrm{Frob}}
\newcommand{\Ind}{\mathrm{Ind}}
\newcommand{\Hom}{\mathrm{Hom}}
\newcommand{\End}{\mathrm{End}}
\newcommand{\Tr}{\mathrm{Tr}}
\newcommand{\Nm}{\mathrm{Nm}}
\newcommand{\ord}{\mathrm{ord}}
\newcommand{\Res}{\mathrm{Res}}
\newcommand{\vol}{\mathrm{vol}}

% 其他
\newcommand{\diam}{\mathrm{diam}}
\newcommand{\Id}{\mathrm{Id}}
\DeclareMathOperator{\im}{im}
\DeclareMathOperator{\coker}{coker}
\DeclareMathOperator{\rank}{rank}
\DeclareMathOperator{\Div}{Div}
\DeclareMathOperator{\divisor}{div}
\DeclareMathOperator{\genus}{genus}
